\par\clearpage
\SourceLeaf{152}\GaussVTheoryEarlyFolio{148}
\[
\begin{aligned}
-K'&=h^{1,1}\frac{\mathrm{d}P^{1,1}}{\mathrm{d}u}+h^{2,1}\frac{\mathrm{d}P^{2,1}}{\mathrm{d}u}+h^{3,1}\frac{\mathrm{d}P^{3,1}}{\mathrm{d}u}+\text{u. s. w.}\\
-l'&=h^{1,1}\frac{P^{1,1}}{\sin u}+h^{2,1}\frac{P^{2,1}}{\sin u}+h^{3,1}\frac{P^{3,1}}{\sin u}+\text{u. s. w.}\\
M'&=2h^{1,1}P^{1,1}+3h^{2,1}P^{2,1}+4h^{3,1}P^{3,1}+\text{u. s. w.}
\end{aligned}
\]
zur Bestimmung der Coëfficienten $h^{1,1}$, $h^{2,1}$, $h^{3,1}$ u. s. w.

Ferner dienen zur Bestimmung der Coëfficienten $g^{2,2}$, $g^{3,2}$, $g^{4,2}$ u. s. w. die Gleichungen
\[
\begin{aligned}
-k''&=g^{2,2}\frac{\mathrm{d}P^{2,2}}{\mathrm{d}u}+g^{3,2}\frac{\mathrm{d}P^{3,2}}{\mathrm{d}u}+g^{4,2}\frac{\mathrm{d}P^{4,2}}{\mathrm{d}u}+\text{u. s. w.}\\
L''&=2g^{2,2}\frac{P^{2,2}}{\sin u}+2g^{3,2}\frac{P^{3,2}}{\sin u}+2g^{4,2}\frac{P^{4,2}}{\sin u}+\text{u. s. w.}\\
m''&=3g^{2,2}P^{2,2}+4g^{3,2}P^{3,2}+5g^{4,2}P^{4,2}+\text{u. s. w.}
\end{aligned}
\]
und auf ähnliche Weise ergeben sich die Coëfficienten der folgenden höhern Ordnungen.

\Article{24}
Der Vorzug dieses Verfahrens vor dem im 22. Art. angegebenen besteht hauptsächlich darin, dass die unbekannten Grössen in Gruppen zerfallen, die jede für sich bestimmt werden, wodurch die Rechnung eine ausserordentliche Erleichterung erhält, während bei dem andern Verfahren die Vermengung sämmtlicher Unbekannten unter einander die Scheidung überaus beschwerlich macht. Dagegen hat jenes Verfahren den Nachtheil, dass es seine Grundlagen gar nicht in unmittelbaren Beobachtungen findet, sondern sie aus graphischen Darstellungen entlehnen muss, welche in den Gegenden, wo Beobachtungen vorhanden sind, diese doch nur roh darstellen können, in solchen Gegenden aber, wo es weit und breit ganz an Beobachtungen fehlt, nur vermuthungsweise, gewissermaassen willkürlich ergänzt sind, und sich daher sehr weit von der Wahrheit entfernen können. Indessen bleibt keine Wahl, als entweder alle Versuche so lange auszusetzen, bis viel vollständigere und zuverlässigere Data bereit sein werden, oder mit den jetzt noch so höchst precären Mitteln einen ersten Versuch zu wagen, von dem man wenig mehr als eine rohe Annäherung erwarten darf. Einen sichern Maassstab für den Werth des Erfolges gibt jedenfalls hinterdrein die scharfe Vergleichung der Resultate mit wirklichen Beobachtungen aus allen Thei-
\par\clearpage
\SourceLeaf{153}\GaussVTheoryEarlyFolio{149}
\noindent len der Erde; und wenn solche Prüfung dahin ausfällt, dass der erste Versuch nicht ganz misslungen ist, so wird dieser eine kräftige Hülfe darbieten, um künftige neue Versuche, auf dem einen oder auf dem andern Wege, zweckmässig vorzubereiten.

\Article{25}
Schon vor vielen Jahren hatte ich zu wiederholten malen angefangen, mich solchen Versuchen zu unterziehen, von denen ich aber immer wieder abzustehen genöthigt war, weil die zu Gebote stehenden Data sich als gar zu dürftig auswiesen. Gleichwohl würde ich schon früher einen Versuch zu Ende zu führen geneigt gewesen sein, wenn der mehrmals von mir ausgesprochene Wunsch in Erfüllung gegangen wäre. dass die reinen horizontalen Intensitäten in einer allgemeinen Karte dargestellt werden möchten, für deren Mangel die Verbindung der vorhandenen unvollkommenen Generalkarten für die Inclination und ganze Intensität keinen Ersatz geben konnte.

Die Erscheinung der \textsc{Sabine}schen Karte für die ganze Intensität (im siebenten \textit{Report of the British association for the advancement of science}) hat mich jetzt zur Unternehmung und Vollendung eines neuen Versuchs angeregt, der übrigens nur aus dem im vorhergehenden Artikel angegebenen Gesichtspunkte angesehen werden soll.

Die der Rechnung unterzulegenden Data wurden aus der erwähnten Karte für die Intensität, der \textsc{Barlow}schen für die Declination (\textit{Philosophical Transactions} 1833), und der von \textsc{Horner} entworfenen für die Inclination (Physikalisches Wörterbuch Band 6.) entnommen, und zwar für je zwölf Punkte auf sieben Parallelkreisen. Die Lücken, welche jene Karten in weiten Strecken übrig lassen, konnten meistens nur auf höchst precäre Art ergänzt werden.

Im Laufe der Rechnung ergab sich bald, dass dieselbe wenigstens bis zu den Grössen der vierten Ordnung ausgedehnt werden müsse, wonach die Anzahl der zu bestimmenden Coëfficienten auf 24 steigt. Aller Wahrscheinlichkeit nach werden auch die Glieder der fünften Ordnung noch ansehnlich genug sein; allein bei einem ersten Versuche bleiben die Werthe von $k$, $m$, $k'$ u. s. w. noch viel zu sehr mit dem Einfluss der vielen unzuverlässigen Daten behaftet, die jener seiner Natur nach einschliessen muss, als dass es verstattet sein könnte, in das Eliminationsgeschäft eine noch grössere Anzahl von unbekannten Grössen aufzunehmen.

\par\clearpage
\SourceLeaf{154}\GaussVTheoryEarlyFolio{150}
Es muss noch bemerkt werden, dass die Intensitäten in \textsc{Sabine}'s Karte dieselbe willkürliche Einheit haben, in welcher sie gewöhnlich bisher ange\-geben zu werden pflegen, und wonach in London die ganze Intensität $=1{,}372$ gesetzt wird. Diese Einheit ist hier bei der Berechnung der Coëfficienten, eben so wie bei der weiter unten zu erklärenden Hülfstafel, dahin abgeändert, dass alle Zahlen tausendmal grösser werden, wobei also die Intensität für London $=1372$ zum Grunde liegt. Übrigens kann offenbar für die Intensität eine jede beliebige Einheit gebraucht werden, insofern man auch die Einheit für $\mu$ als willkürlich betrachten, und diese immer jener gemäss annehmen kann. Will man weitere Folgerungen daran knüpfen, für welche $\mu$ auf ein absolutes Maass gebracht sein muss, so brauchen nur sämmtliche Coëfficienten mit demselben Factor multiplicirt zu werden, welcher zur Reduction der nach jener Einheit ausgedrückten Intensitätszahlen auf absolutes Maass erforderlich ist.

\Article{26}
Die aus der ersten Rechnung, wobei die Längen $\lambda$ von Greenwich östlich gezählt sind, erhaltenen Zahlwerthe der 24 Coëfficienten sind folgende:
\[
\begin{array}{r|r}
g^{1,0}=+925{,}782&g^{2,2}=+0{,}493\\
g^{2,0}=-22{,}059&g^{3,2}=-73{,}193\\
g^{3,0}=-18{,}868&g^{4,2}=-45{,}791\\
g^{4,0}=-108{,}855&h^{2,2}=-39{,}010\\
g^{1,1}=+89{,}024&h^{3,2}=-22{,}766\\
g^{2,1}=-144{,}913&h^{4,2}=+42{,}573\\
g^{3,1}=+122{,}936&g^{3,3}=+1{,}396\\
g^{4,1}=-152{,}589&g^{4,3}=+19{,}774\\
h^{1,1}=-178{,}744&h^{3,3}=-18{,}750\\
h^{2,1}=-6{,}030&h^{4,3}=-0{,}178\\
h^{3,1}=+47{,}794&g^{4,4}=+4{,}127\\
h^{4,1}=+64{,}112&h^{4,4}=+3{,}175
\end{array}
\]
Diese Zahlen, welche man als die \textit{Elemente der Theorie des Erdmagnetismus} betrachten kann, sind hier genau so angesetzt, und als Grundlage der nachher zu beschreibenden Hülfstafel angewandt, wie die Rechnung sie gegeben hat, ohne die Decimalbrüche wegzulassen. Für jeden Rechnungskundigen ist die Be-
\par\clearpage
\SourceLeaf{155}\GaussVTheoryEarlyFolio{151}
\noindent merkung überflüssig, dass diese Bruchtheile an sich keinen Werth haben, da wir noch weit davon entfernt sind, nur die ganzen Einer mit Zuverlässigkeit ausmitteln zu können: allein es ist von Wichtigkeit, dass die Beobachtungen mit einem und demselben bestimmten System von Elementen scharf verglichen werden, und da war kein Grund vorhanden, an dem, was die Rechnung ergeben hatte, etwas zu verändern, weil durch Weglassung der Decimalbrüche für die Bequemlichkeit der Vergleichungsrechnungen gar nichts gewonnen worden sein würde.

\Article{27}
Der entwickelte Ausdruck für $V$ nach obigen Zahlen ist folgender, wobei der Abkürzung wegen $e$ für $\cos u$ und $f$ für $\sin u$ geschrieben ist.
\[
\begin{aligned}
\frac VR={}&-1{,}977+937{,}103e+71{,}245ee-18{,}868e^3-108{,}855e^4\\
&+(64{,}437-79{,}518e+122{,}936ee+152{,}589e^3)f\cos\lambda\\
&+(-188{,}303-33{,}507e+47{,}794ee+64{,}112e^3)f\sin\lambda\\
&+(7{,}035-73{,}193e-45{,}791ee)ff\cos2\lambda\\
&+(-45{,}092-22{,}766e-42{,}573ee)ff\sin2\lambda\\
&+(1{,}396+19{,}774e)f^3\cos3\lambda\\
&+(-18{,}750-0{,}178e)f^3\sin3\lambda\\
&+4{,}127f^4\cos4\lambda\\
&+3{,}175f^4\sin4\lambda
\end{aligned}
\]

Es mögen ferner die vollständig entwickelten Ausdrücke für die drei Componenten der magnetischen Kraft hier Platz finden.
\[
\begin{aligned}
X={}&(937{,}103+142{,}490e-56{,}603ee-435{,}420e^3)f\\
&+(-79{,}518+181{,}435e-298{,}732ee\\
&\qquad-368{,}808e^3+610{,}357e^4)\cos\lambda\\
&+(-33{,}507+283{,}892e+259{,}349ee\\
&\qquad-143{,}383e^3-256{,}448e^4)\sin\lambda\\
&+(-73{,}193-105{,}652e+219{,}579ee+183{,}164e^3)f\cos2\lambda\\
&+(-22{,}766+175{,}330e+68{,}098ee-170{,}292e^3)f\sin2\lambda\\
&+(19{,}774-4{,}188e-79{,}096ee)ff\cos3\lambda\\
&+(-0{,}178+56{,}250e+0{,}716ee)ff\sin3\lambda\\
&-16{,}508ef^3\cos4\lambda\\
&-12{,}701ef^3\sin4\lambda
\end{aligned}
\]
\par\clearpage
\SourceLeaf{156}\GaussVTheoryEarlyFolio{152}
\[
\begin{aligned}
Y={}&(188{,}303+33{,}507e-47{,}794ee-64{,}112e^3)\cos\lambda\\
&+(64{,}437-79{,}518e+122{,}936ee-152{,}589e^3)\sin\lambda\\
&+(90{,}184+45{,}532e-85{,}146ee)f\cos2\lambda\\
&+(14{,}070-146{,}386e-91{,}582ee)f\sin2\lambda\\
&+(56{,}250+0{,}534e)ff\cos3\lambda\\
&+(4{,}188+59{,}322e)ff\sin3\lambda\\
&-12{,}701f^3\cos4\lambda\\
&+16{,}508f^3\sin4\lambda
\end{aligned}
\]
\[
\begin{aligned}
Z={}&-24{,}593+1896{,}847e+400{,}343ee-75{,}471e^3-544{,}275e^4\\
&+(79{,}700-107{,}763e+491{,}744ee-762{,}946e^3)f\cos\lambda\\
&+(-395{,}724-155{,}473e+191{,}176ee+320{,}560e^3)f\sin\lambda\\
&+(34{,}187-292{,}772e-228{,}955ee)ff\cos2\lambda\\
&+(-147{,}439-91{,}064e+212{,}865ee)ff\sin2\lambda\\
&+(5{,}584+98{,}870e)f^3\cos3\lambda\\
&+(-75{,}000-0{,}890e)f^3\sin3\lambda\\
&+20{,}635f^4\cos4\lambda\\
&+15{,}876f^4\sin4\lambda
\end{aligned}
\]

Nachdem diese Componenten für einen gegebenen Ort berechnet sind, erhält man die Bestimmungsstücke der magnetischen Kraft in der gewöhnlichen Form auf folgende Art. Es sei $\delta$ die Declination, $i$ die Inclination, $\psi$ die ganze, $\omega$ die horizontale Intensität. Man bestimmt zuerst $\delta$ und $\omega$ vermittelst der Formeln
\[
X=\omega\cos\delta,\qquad Y=\omega\sin\delta
\]
und sodann $i$ und $\psi$ vermittelst der folgenden
\[
\omega=\psi\cos i,\qquad Z=\psi\sin i
\]

\Article{28}
Da die Formeln für $X$, $Y$, $Z$ zusammen 71 Glieder enthalten, so ist die unmittelbare Rechnung nach denselben eine ziemlich beschwerliche Arbeit, und die Wiederholung derselben für eine grosse Anzahl von Örtern würde allerdings desto mehr abschreckendes haben, da man ohne dieselbe Rechnung zweimal zu machen nicht wohl hoffen dürfte, gegen mögliche Rechnungsfehler geschützt zu
